% =====================================================================
%  CHƯƠNG 5. NHẬN DẠNG KÝ TỰ VÀ MÀU NỀN BIỂN SỐ   | Nhật (OCR) + Đức (màu)
%  Ngân sách: mục tiêu 7 trang. Viết high-level, không nhắc mã nguồn/hàm/đường dẫn.
%  Nguồn số nội bộ (không in): ocr_full_progression.csv, ocr_comparison.csv,
%    ocr_vnplate_results.csv (a1/vn_plate), topkek 37 ký tự.
% =====================================================================
\chapter{Nhận Dạng Ký Tự và Màu Nền Biển Số}
\label{ch:ocr}

% =====================================================================
\section{Dữ Liệu OCR}
\label{sec:ocr-dulieu}
% =====================================================================

\subsection{Dữ Liệu Huấn Luyện}

Nguồn chính là bộ topkek\_plate\_ocr~\autocite{topkek2024vnplateocr}, gồm 6.643 crop biển thật
và 5.547 crop sinh tổng hợp. Bố cục
biển được suy từ tỉ lệ khung ảnh (rộng trên cao nhỏ hơn 2,0 coi là biển hai dòng). Tập chia
80/10/10 phân tầng theo bố cục, giữ crop sinh tổng hợp riêng và chỉ dùng khi huấn luyện
(Bảng~\ref{tab:ocr-split}, phân bố ở Hình~\ref{fig:ocr-data-dist}).

\begin{table}[htbp]
\centering
\caption{Phân chia tập dữ liệu OCR theo split và bố cục biển.}

{\footnotesize Nguồn: bộ topkek\_plate\_ocr~\autocite{topkek2024vnplateocr}, chia phân tầng theo tỉ lệ bố cục.}
\label{tab:ocr-split}
\begin{tabular}{lrrr}
\toprule
Split & Biển 1 dòng & Biển 2 dòng & Tổng \\
\midrule
train & 1.496 & 3.817 & 5.313 \\
val   & 187   & 477   & 664   \\
test  & 188   & 478   & 666   \\
Sinh tổng hợp (chỉ train) & \multicolumn{2}{c}{(không tách theo bố cục)} & 5.547 \\
\bottomrule
\end{tabular}
\end{table}

\begin{figure}[htbp]
\centering
\includegraphics[width=0.85\textwidth]{plate_ocr_data_distribution.png}
\caption{Phân bố dữ liệu OCR theo split và bố cục biển.}

\label{fig:ocr-data-dist}
\end{figure}

\subsection{Tập Test Độc Lập}

Tập test của topkek cùng nguồn gốc với tập huấn luyện nên chưa đủ đánh giá tổng quát hoá. Đề
tài xây thêm tập test độc lập từ bộ vn\_plate (Roboflow, CC BY 4.0, 1.005 ảnh) có sẵn bounding
box nhưng không có nhãn chuỗi: cắt biển theo bounding box, chọn ngẫu nhiên 100 biển rồi gán
nhãn thủ công mà không nhìn dự đoán của mô hình, giữ lại 96 biển đọc được (45 biển một dòng,
51 biển hai dòng). Tập này có độ phân giải sát điều kiện triển khai hơn: trung vị 52 px mỗi
dòng ký tự, so với 19 px ở topkek.

% =====================================================================
\section{Kiến Trúc}
\label{sec:ocr-kientruc}
% =====================================================================

Đề tài so sánh mô hình tự huấn luyện với hai engine OCR pretrained làm mốc: RapidOCR (dòng
PP-OCRv3, chạy qua ONNX Runtime)~\autocite{rapidocr} và EasyOCR (CRNN pretrained, bảng ký tự
tiếng Anh)~\autocite{easyocr}. Mô hình riêng là một CRNN gồm phần trích đặc trưng tích chập,
một lớp BiLSTM và đầu ra CTC, ảnh vào ảnh xám 48$\times$128 pixel, tổng khoảng 425 nghìn tham
số, nhỏ nhờ bảng ký tự thu hẹp đúng cấu trúc biển Việt Nam; kiến trúc theo Shi và cộng sự
cùng Graves và cộng sự~\autocite{shi2017crnn,graves2006ctc}. Cả ba dùng chung luồng xử lý:
chỉnh hình học, tách dòng nếu là biển hai dòng, đọc từng dòng, ghép chuỗi và chuẩn hoá theo
quy tắc biển Việt Nam.

\subsection{Nắn Phối Cảnh Thay Cho Xoay Phẳng}

Xoay biển trong mặt phẳng ảnh không xử lý được biến dạng phối cảnh khi camera gắn cao và lệch
bên. Đề tài nắn 4 góc biển, lấy từ đầu ra phân vùng của mô hình phát hiện, về hình chữ nhật
bằng phép biến đổi phối cảnh, tăng tỉ lệ đọc đúng trên A1 từ 44\% lên 76\%
(Bảng~\ref{tab:ocr-perspective}).

\begin{table}[htbp]
\centering
\caption{So sánh hai cách xử lý hình học trên tập A1 (ảnh cổng bãi xe thật).}

{\footnotesize Đếm thủ công trên các biển đọc được bằng mắt trong tập A1.}
\label{tab:ocr-perspective}
\begin{tabular}{lr}
\toprule
Cách xử lý hình học & Đọc đúng \\
\midrule
Xoay phẳng & 8/18 (44\%) \\
\textbf{Nắn phối cảnh 4 điểm} & \textbf{13/17 (76\%)} \\
\bottomrule
\end{tabular}
\end{table}

% =====================================================================
\section{Huấn Luyện và Xử Lý Dữ Liệu}
\label{sec:ocr-caitien}
% =====================================================================

Mỗi mẫu huấn luyện là một dòng ký tự; biển hai dòng được tách thành hai mẫu. Bảng ký tự ban
đầu gồm 36 ký tự (0 đến 9, A đến Z), mở rộng thêm ``Đ'' ở Mục~\ref{sec:ocr-charset37}; siêu
tham số chi tiết ở Bảng~\ref{tab:app-ocr}.

\subsection{Xử Lý Dữ Liệu Sinh Tổng Hợp}

Ảnh sinh tổng hợp chiếm khoảng một nửa dữ liệu huấn luyện nhưng là biển render sắc nét, không
nhiễu, không nghiêng, trung vị 45 đến 50 px mỗi dòng so với 19 px của crop thật. Đề tài so
sánh ba cách dùng ảnh này, các yếu tố khác giữ nguyên: giữ nguyên, bỏ hẳn, và hạ cấp cho khớp
phân phối ảnh thật (thu nhỏ theo phân phối px/dòng thực nghiệm rồi thêm nhiễu và nén). Cách
hạ cấp tốt nhất trên cả hai tập test, bỏ hẳn là kém nhất (Bảng~\ref{tab:ocr-v0v2}); lọc bỏ
crop dưới 25 px/dòng cũng cho kết quả kém hơn. Số lượng dữ liệu vì vậy quan trọng hơn độ đồng
nhất phân phối.

\begin{table}[htbp]
\centering
\caption{Accuracy theo cách xử lý dữ liệu sinh tổng hợp (bảng 36 ký tự).}

\label{tab:ocr-v0v2}
\begin{tabular}{lrr}
\toprule
Cách xử lý & topkek test & vn\_plate test \\
\midrule
Giữ nguyên & 51,1\% & 62,5\% \\
Bỏ hẳn ảnh sinh tổng hợp & 49,5\% & 51,0\% \\
\textbf{Hạ cấp khớp phân phối} & \textbf{54,2\%} & \textbf{64,6\%} \\
\bottomrule
\end{tabular}
\end{table}

\subsection{Ranh Giới Nhãn Biển Hai Dòng}

Nhãn biển hai dòng là một chuỗi ghép không chỉ rõ ranh giới dòng; tách sai khiến mô hình học
bỏ ký tự cuối dòng trên. Tách nhãn theo dấu cách trong nhãn gốc sửa được khoảng 10,6\% nhãn
biển hai dòng bị tách sai và gần như xoá chênh lệch giữa biển một dòng và hai dòng.

\subsection{Kiểm Soát Nhiễu Ngẫu Nhiên Giữa Các Lần Huấn Luyện}

Hai lần huấn luyện cùng cấu hình (bảng 36 ký tự, sau khi sửa nhãn hai dòng), chỉ khác seed 42
và 123, cho accuracy trên vn\_plate 87,5\% và 92,7\%, lệch 5,2 điểm phần trăm. Với tập test
nhỏ, mức lệch này đủ làm sai kết luận rút ra từ một lần huấn luyện, nên đề tài cố định seed và
chọn mô hình theo CER trên tập kiểm định thay vì theo điểm cao nhất trên tập test. Hình~\ref{fig:ocr-loss-curves} trình bày đường mất mát và
CER kiểm định theo epoch.

\begin{figure}[htbp]
\centering
\includegraphics[width=0.90\textwidth]{plate_ocr_progression.png}
\caption{Đường mất mát huấn luyện và CER/accuracy kiểm định theo epoch.}

\label{fig:ocr-loss-curves}
\end{figure}

% =====================================================================
\section{Mở Rộng Bảng Ký Tự Lên 37 Ký Tự và Model Triển Khai}
\label{sec:ocr-charset37}
% =====================================================================

Để đọc seri MĐ và TĐ của xe máy điện, đề tài mở bảng ký tự lên 37 ký tự và giữ ``Đ'' trong
toàn bộ khâu chuẩn hoá nhãn và hậu xử lý, vốn từng vô tình loại bỏ ký tự này. Mô hình triển
khai được đo trên ba tập test có độ phân giải tăng dần (Bảng~\ref{tab:ocr-deployed}).

\begin{table}[htbp]
\centering
\caption{Kết quả model triển khai CRNN 37 ký tự trên ba tập test (accuracy mức cả biển).}

\label{tab:ocr-deployed}
\begin{tabular}{lrrr}
\toprule
& topkek (n = 666) & vn\_plate (n = 96) & a1 (n = 172) \\
\midrule
Accuracy toàn biển & 56,6\% & 80,2\% & 91,9\% \\
Biển 1 dòng & 59,6\% & 75,6\% & 91,5\% \\
Biển 2 dòng & 55,4\% & 84,3\% & 92,3\% \\
CER & 0,183 & 0,031 & 0,015 \\
\bottomrule
\end{tabular}
\end{table}

Tập a1 (172 crop cắt từ ảnh camera cổng, đã nắn phối cảnh) gần nhất với điều kiện chạy thật;
trên tập này mô hình đạt \textbf{91,9\%} accuracy mức cả biển (khoảng tin cậy Wilson 95\% từ
86,8\% đến 95,1\%) với CER 0,015, nên là con số dùng để đối chiếu tiêu chí đọc biển trên
90\% của đề cương. Trên cả ba tập, biển hai dòng không kém biển một
dòng quá 4,2 điểm phần trăm, cho thấy bước tách dòng xử lý tốt cả hai bố cục.

Đóng góp của ký tự ``Đ'' là mở rộng độ phủ chứ không phải tăng accuracy tổng thể: trên 17
biển có ``Đ'' trong tập test, mô hình 37 ký tự đọc đúng ``Đ'' ở 11/17 biển (khớp toàn chuỗi
ở 8/17), trong khi mô hình 36 ký tự cũ không đọc được biển nào (0/17). Bản ONNX cho dự đoán
trùng khớp bản huấn luyện trên toàn bộ 666 biển kiểm tra.

% =====================================================================
\section{Kết Quả So Sánh và Phân Tích Độ Phân Giải}
\label{sec:ocr-ketqua}
% =====================================================================

\subsection{So Sánh Với Hai Engine Pretrained}

Trên tập test topkek (666 crop thật), mô hình CRNN riêng vượt cả hai baseline trên mọi tiêu
chí cùng lúc (Bảng~\ref{tab:ocr-compare}): chính xác hơn nhiều lần, nhỏ hơn RapidOCR khoảng 6
lần và EasyOCR khoảng 58 lần, và nhanh hơn cả hai trên CPU. Nguyên nhân là phạm vi bài toán:
hai mô hình pretrained mang bảng ký tự tổng quát, còn CRNN riêng chỉ cần phân biệt bộ ký tự
cố định theo cấu trúc biển Việt Nam (Hình~\ref{fig:ocr-compare}).

\begin{table}[htbp]
\centering
\caption{So sánh ba phương pháp OCR trên tập test topkek (666 biển thật).}

{\footnotesize CRNN riêng đo ở bảng 36 ký tự, trước bước mở rộng ``Đ''; kích thước tính cả mô hình phát hiện văn bản của EasyOCR; thời gian đo trên CPU máy huấn luyện, cho một lần đọc.}
\label{tab:ocr-compare}
\begin{tabular}{lrrrr}
\toprule
Phương pháp & Accuracy & CER & Kích thước & CPU inference \\
\midrule
RapidOCR~\autocite{rapidocr} & 34,2\% & 34,2\% & 10,7 MB & 7,20 ms \\
EasyOCR~\autocite{easyocr}   & 11,7\% & 51,9\% & 98,3 MB & 16,65 ms \\
\textbf{CRNN riêng}         & \textbf{58,9\%} & \textbf{17,4\%} & \textbf{1,7 MB} & \textbf{1,69 ms} \\
\bottomrule
\end{tabular}
\end{table}

\begin{figure}[htbp]
\centering
\includegraphics[width=0.82\textwidth]{plate_ocr_accuracy_comparison.png}
\caption{Accuracy toàn biển và CER của ba phương pháp trên tập test topkek.}

\label{fig:ocr-compare}
\end{figure}

\subsection{Phân Tích Theo Dải Độ Phân Giải}

Cùng một mô hình cho hai con số cách nhau 24 điểm phần trăm trên hai tập test (56,6\% so với
80,2\%). Tách theo dải độ phân giải (Bảng~\ref{tab:ocr-resolution-band}) cho thấy nguyên nhân
là tỉ lệ ảnh khó giữa hai tập rất khác nhau, không phải mô hình thiếu ổn định: ở dải đủ nét
(từ 40 px/dòng trở lên) mô hình đạt 93,9\% trên topkek, vượt mục tiêu 90\%, nhưng dải đó chỉ
chiếm 7,4\% tập topkek trong khi hơn nửa tập dưới 20 px, gần như không đọc được bằng mắt. Tập
vn\_plate có phân bố ngược (87,5\% ảnh từ 40 px trở lên) nên trung bình cao hơn dù dùng chung
mô hình. Vì vậy accuracy toàn tập chỉ có ý nghĩa khi đọc kèm phân bố độ phân giải. Một số mẫu dự đoán ở Hình~\ref{fig:ocr-samples}.

\begin{table}[htbp]
\centering
\caption{Accuracy theo dải độ phân giải (bảng 36 ký tự). Tập vn\_plate không có ảnh dưới 20 px/dòng.}

\label{tab:ocr-resolution-band}
\begin{tabular}{lrrrr}
\toprule
Dải độ phân giải & Tỉ lệ trong topkek & Accuracy (topkek) & Tỉ lệ trong vn\_plate & Accuracy (vn\_plate) \\
\midrule
Dưới 20 px/dòng & 52,3\% & 46,6\% & 0\% & n/a \\
20 đến 30 px    & 25,8\% & 61,6\% & 1,0\% & 100\% \\
30 đến 40 px    & 14,6\% & 80,4\% & 11,5\% & 100\% \\
\textbf{Từ 40 px trở lên} & 7,4\% & \textbf{93,9\%} & 87,5\% & 85,7\% \\
\bottomrule
\end{tabular}
\end{table}

\begin{figure}[htbp]
\centering
\includegraphics[width=0.82\textwidth]{plate_ocr_sample_predictions.png}
\caption{Mẫu dự đoán của mô hình CRNN 37 ký tự trên tập test.}

\label{fig:ocr-samples}
\end{figure}

% =====================================================================
\section{Hậu Xử Lý Theo Quy Định Biển Số Việt Nam}
\label{sec:ocr-hauxy}
% =====================================================================

Sau khi ghép chuỗi, kết quả được chuẩn hoá theo vị trí kỳ vọng: vị trí số quy về chữ số (O
thành 0, I thành 1, B thành 8, S thành 5, Z thành 2), vị trí seri quy về chữ cái, rồi kiểm
tra khớp định dạng. Seri có thể gồm một hoặc hai chữ cái: theo Thông tư
79/2024/TT-BCA~\autocite{tt792024bca} (hiệu lực 01/01/2025), biển xe máy thống nhất seri hai
chữ cái kèm 5 chữ số, biển seri cũ một chữ một số chỉ dùng đến hết 31/12/2025, và biển seri
hai chữ cái đã chiếm 15\% đến 19\% tập topkek. Biểu thức kiểm tra nhận cả hai trường hợp. Với
biển hai dòng, độ dài dòng trên cho biết chỗ ngắt khi hiển thị nên không cần suy đoán; chỉ
biển một dòng mới phải suy đoán theo độ dài phần số.

% =====================================================================
\section{Nhận Diện Màu Nền Biển Số}
\label{sec:mau-bien}
% =====================================================================

\subsection{Vai Trò và Phương Pháp}

Màu nền biển xác định nhóm chủ xe theo pháp luật (Bảng~\ref{tab:maubienso}) và được hệ thống
dùng để phân nhóm xe, áp biểu phí (Chương~\ref{ch:hethong}); phân bố màu trên dữ liệu thật ở
Bảng~\ref{tab:hsv-color}.

Đề tài dùng không gian màu HSV vì kênh Hue và Saturation tách rõ màu sắc khỏi độ sáng, giảm
ảnh hưởng biến thiên ánh sáng: trắng có Saturation thấp; vàng, xanh dương và đỏ nằm ở các dải
Hue riêng biệt. Thuật toán đếm pixel nền thuộc từng cụm màu rồi chọn nhãn chiếm đa số, kèm
điểm tin cậy theo tỉ lệ pixel nền được nhận diện.

\subsection{Khử Nhiễu Viền Bằng Otsu và Phát Hiện Contour}

Crop biển thường chứa viền thừa lẫn màu sơn xe hoặc viền đen, gây nhầm màu (ví dụ một biển
trắng từng bị phân loại thành xanh do pixel ở dải viền). Đề tài bổ sung bước tách vùng nền
trước khi đếm màu: tăng tương phản cục bộ bằng CLAHE, phân tách vùng sáng và tối bằng ngưỡng
Otsu, lấy vùng sáng liền khối lớn nhất làm nền biển rồi cắt động theo vùng đó. Cơ chế an toàn
giữ nguyên crop gốc nếu vùng nền tìm được quá nhỏ (dưới 30\% diện tích), tránh cắt sai.

\subsection{Kết Quả và Giới Hạn 2/4 Màu}

Trên tập nhãn màu hiện có gồm 178 crop hai lớp trắng và đỏ, mô đun phân loại đúng toàn bộ
(178/178), gồm cả biển đỏ chữ sáng trên nền tối vốn từng bị nhầm thành trắng. Tập này chưa có
nhãn màu vàng và xanh nên mới đánh giá định lượng được 2 trong 4 màu; biển vàng dưới đèn
huỳnh quang có nguy cơ bị nhầm thành trắng do saturation giảm.

% =====================================================================
\section{Thảo Luận}
\label{sec:ocr-thao-luan}
% =====================================================================

Độ chính xác OCR chịu tác động của ba yếu tố độc lập, mỗi yếu tố một cách xử lý: chất lượng
ảnh (thêm dữ liệu đúng phân phối, không lọc bớt), góc chụp chéo (nắn phối cảnh) và chất lượng
nhãn (tách đúng ranh giới hai dòng).

Về triển khai, mỗi dòng ký tự cần tối thiểu khoảng 40 px. Con số này phụ thuộc vị trí và góc
camera hơn là số megapixel (ảnh gốc A1 rộng tới 4.032 px nhưng biển chỉ khoảng 31 px/dòng vì
xe ở xa), nên cần đặt camera gần hoặc thu hẹp góc nhìn thay vì mua camera phân giải cao hơn.
Tập vn\_plate chỉ có 96 biển nên số liệu trên tập này nên đọc như xu hướng.
